\section{The reduced optical action}
\label{sec:v109-action}
Fix an actual collision itinerary with $m\geq1$ flights. Its visited
centres are $O_0,\ldots,O_m$, its collision points are
$x_i=O_i+R n(\alpha_i)$, and $n(\alpha)=(\cos\alpha,\sin\alpha)$.
Write $\tau_i=|x_{i+1}-x_i|$. All variations in this section are made
on this actual itinerary before a physical boundary is reached.
They neither continue an event through a forbidden collision nor
change an occupation label.

\begin{lemma}[The original section stays away from endpoint grazing]
\label{lem:v109-section}
On the closure of the original initial and terminal section tests,
\begin{equation}\label{eq:v109-endpoint-cap}
 |p_0|,|p_m|\leq\frac9{10},\qquad
 c_0,c_m\geq c_\circ:=\frac{\sqrt{19}}{10}.
\end{equation}
\end{lemma}
\begin{proof}
The only section centres with tangential coordinate exceeding $1/2$
in absolute value have coordinate $2q_R/(1+q_R^2)$.
This function is increasing on $(0,1)$, and
\[
 \frac{2(53/100)}{1+(53/100)^2}+\frac1{200}<\frac9{10}.
\]
The other squares and the rectangle $U$ satisfy the same strict
bound. Passing to their closures proves the assertion.
\end{proof}

Let $S(\alpha_0,\ldots,\alpha_m)=\sum_{i=0}^{m-1}\tau_i$ be the
broken-path action. At a regular specular trajectory its interior
variables are stationary. Eliminating them locally gives the actual
endpoint action $F(q)$, where $q=(\alpha_0,\alpha_m)$.

\begin{lemma}[Second variation and its Schur complement]
\label{lem:v109-hessian}
For a real variation $\xi=(\xi_0,\ldots,\xi_m)$ at an actual
regular trajectory,
\begin{equation}\label{eq:v109-second-variation}
 \begin{split}
 D^2S[\xi,\xi]={}&R\left(c_0\xi_0^2+c_m\xi_m^2
                   +2\sum_{i=1}^{m-1}c_i\xi_i^2\right)\\
 &+R^2\sum_{i=0}^{m-1}
                 \frac{(c_i\xi_i+c_{i+1}\xi_{i+1})^2}{\tau_i}.
 \end{split}
\end{equation}
The reduced Hessian $G=D_q^2F$ satisfies
\begin{equation}\label{eq:v109-hessian-bounds}
 R\operatorname{diag}(c_0,c_m)\leq G,\qquad
 \lambda_\circ I_2\leq G\leq\Lambda_\circ I_2,
 \quad \lambda_\circ=\frac{9\sqrt{19}}{200},\quad
 \Lambda_\circ=\frac{47}{6}.
\end{equation}
These inequalities are uniform in the radius, the collision count,
and the original exact labels.
\end{lemma}
\begin{proof}
For a flight with unit direction $v$, the Hessian of its length in
its endpoints is the squared perpendicular displacement divided by
the flight length. The acceleration of $R n(\alpha)$ is $-R n(\alpha)$.
At the departure and arrival points, respectively,
$v\cdot n_i=c_i$ and $v\cdot n_{i+1}=-c_{i+1}$.
Also $(R n_i')\cdot v^\perp=Rc_i$ and
$(R n_{i+1}')\cdot v^\perp=-Rc_{i+1}$.
Thus one flight contributes
$R(c_i\xi_i^2+c_{i+1}\xi_{i+1}^2)$ and the corresponding
square in \eqref{eq:v109-second-variation}. Summing proves that formula.

Let $H=D^2S$, $E=\{0,m\}$ and $I=\{1,\ldots,m-1\}$.
For $m>1$, the interior block is positive definite. The implicit
function theorem and differentiation of the stationary action give
\begin{equation}\label{eq:v109-schur}
 G=H_{EE}-H_{EI}H_{II}^{-1}H_{IE},\qquad
 \eta^{\mathsf T}G\eta=
    \min_{\zeta}(\eta,\zeta)^{\mathsf T}H(\eta,\zeta).
\end{equation}
The order in the last expression places endpoint coordinates first.
The diagonal endpoint terms in \eqref{eq:v109-second-variation}
prove the first lower bound. Choosing the interior variation to be
zero gives $G\leq H_{EE}$. The latter block is bounded above by
$\{R_++2R_+^2/\tau_-\}I_2=(47/6)I_2$; the same bound covers the
single-flight case, where there is no Schur complement.
Lemma~\ref{lem:v109-section} proves the remaining lower bound.
\end{proof}

\subsection{A positive optical denominator}
Put
\begin{equation}\label{eq:v109-optical-matrices}
 \mathsf S(s)=\begin{pmatrix}1&s\\0&1\end{pmatrix},\qquad
 \mathsf C(\kappa)=\begin{pmatrix}1&0\\\kappa&1\end{pmatrix},
 \qquad \kappa_i=\frac{2}{Rc_i}.
\end{equation}
Let $\mathsf M_m$ be the product from the initial collision to the
terminal collision, with the rightmost factor acting first:
\begin{equation}\label{eq:v109-optical-product}
 \mathsf M_m=
 \mathsf S(\tau_{m-1})\mathsf C(\kappa_{m-1})\cdots
 \mathsf S(\tau_1)\mathsf C(\kappa_1)\mathsf S(\tau_0),
 \qquad \mathfrak b_m=(\mathsf M_m)_{12}.
\end{equation}
For $m=1$ the product is just $\mathsf S(\tau_0)$.

\begin{lemma}[Exact denominator and a marked incidence bound]
\label{lem:v109-denominator}
The actual endpoint-coordinate density is
\begin{equation}\label{eq:v109-omega}
 \omega=|\partial_{p_0}\alpha_m|^{-1}
        =\frac{|G_{0m}|}{R}=\frac{Rc_0c_m}{\mathfrak b_m}.
\end{equation}
With $\varrho=47/59$, one has
\begin{equation}\label{eq:v109-optical-bounds}
 \omega\leq\frac{47}{6}\varrho^{m-1},\qquad
 \frac{\omega}{c_j}\leq\frac{2209}{72}\varrho^{m-2}
          \quad(1\leq j\leq m-1).
\end{equation}
The second inequality is asserted when $m\geq2$.
\end{lemma}
\begin{proof}
The first equality with $G$ follows by differentiating
$p_0=-\partial_{\alpha_0}F/R$ with $\alpha_0$ fixed.
For the denominator, the interior block of $H$ is
$R^2 C_I J C_I$, where $C_I=\operatorname{diag}(c_1,\ldots,c_{m-1})$
and $J$ is tridiagonal with
\[
 J_{ii}=\kappa_i+\tau_{i-1}^{-1}+\tau_i^{-1},\qquad
 J_{i,i+1}=\tau_i^{-1}.
\]
The corner cofactor in \eqref{eq:v109-schur} gives
$|G_{0m}|=R^2c_0c_m/(\prod_i\tau_i\det J)$.
Expansion of successive tridiagonal determinants gives
$\prod_i\tau_i\det J=\mathfrak b_m$. This identity also follows
by multiplying the second column in \eqref{eq:v109-optical-product};
both recursions start with $\mathfrak b_1=\tau_0$.
The empty determinant convention handles $m=1$.

All entries in the optical factors are nonnegative. Multiplication
of a positive second column $(b,d)^{\mathsf T}$ by
$\mathsf S(\tau_i)\mathsf C(\kappa_i)$ changes its top entry to
$(1+\tau_i\kappa_i)b+\tau_i d$.
Since $1+\tau_i\kappa_i\geq1+2\tau_-/R_+=59/47$,
\begin{equation}\label{eq:v109-denominator-lower}
 \mathfrak b_m\geq\tau_-\varrho^{-(m-1)}.
\end{equation}
For a fixed interior collision $j$, split the product immediately
before and after $\mathsf C(\kappa_j)$. Its $(1,2)$ entry is an
expression with nonnegative coefficients in $\kappa_j$, whose linear
coefficient is $\mathfrak b_{\rm right}\mathfrak b_{\rm left}$.
The two subproducts have respectively $m-j$ and $j$ flights. Hence
\[
 \mathfrak b_m\geq
   \kappa_j\tau_-^2\varrho^{-(m-2)}.
\]
Substitution into \eqref{eq:v109-omega}, and $c_0,c_m\leq1$, proves
\eqref{eq:v109-optical-bounds}.
\end{proof}

\begin{corollary}[The full inherited incidence insertion]
\label{cor:v109-marked}
For the unchanged insertion
$I_m=1+\sum_{i=1}^{m-1}q_0^{d_i}/c_i$, with $q_0=47/53$,
\begin{equation}\label{eq:v109-marked-bound}
 \omega I_m\leq500\varrho^{(m-2)_+}.
\end{equation}
This is an estimate on each actual physical trajectory; it involves
no lower cutoff on any interior incidence.
\end{corollary}
\begin{proof}
For $m\geq2$,
$\sum_{i=1}^{m-1}q_0^{d_i}\leq2q_0/(1-q_0)=47/3$.
Use \eqref{eq:v109-optical-bounds} and
$(47/6)\varrho+(2209/72)(47/3)<500$.
For $m=1$ the sum is empty and the first optical bound suffices.
\end{proof}
No estimate for a sum over collision words is obtained by simply
summing the right side of \eqref{eq:v109-marked-bound}. In particular,
this per-trajectory exponential estimate is not being substituted for
a collision-uniform local limit theorem.

\section{Completion of the original roofwise endpoint traces}
\label{sec:v109-traces}
We use the original incident-component pairing sheets and their
constant-angle coordinates $(u,\theta)=(F,\arg\nabla F)$.
On one of their positive-length maximal regular segments, write
$a<u<b$, $g=\nabla F=z e(\theta)$ and
$e^\perp(\theta)=(-\sin\theta,\cos\theta)$.
The physical section and all original exact-event tests are retained.

\begin{lemma}[Finite Gauss-coordinate length]
\label{lem:v109-gauss}
Along a fixed-$\theta$ segment,
\begin{equation}\label{eq:v109-gauss}
 \frac{dq}{dz}=G^{-1}e,\qquad
 \frac{du}{dz}=z e^{\mathsf T}G^{-1}e,\qquad
 2\lambda_\circ\leq\frac{d(z^2)}{du}\leq2\Lambda_\circ.
\end{equation}
Consequently $z$ is strictly increasing, both endpoint limits of
$q$ exist, and
\begin{equation}\label{eq:v109-integrable-inverse-gradient}
 \int_a^b\frac{du}{z(u,\theta)}
 \leq\lambda_\circ^{-1}(z_b-z_a)
 \leq\lambda_\circ^{-1}\min\{\sqrt2 R_+,
                                      \sqrt{2\Lambda_\circ(b-a)}\}.
\end{equation}
A critical endpoint of a positive-length segment can only be its
incoming endpoint.
\end{lemma}
\begin{proof}
Differentiate $\nabla F(q)=z e$ and then $F(q)=u$.
The bounds in \eqref{eq:v109-hessian-bounds} give
\eqref{eq:v109-gauss}. They also give $|dq/dz|\leq\lambda_\circ^{-1}$.
Since $0<z\leq\sqrt2 R_+$, the two endpoint limits of $q$ exist
in the chosen angular lifts. Integrating $du/z=e^{\mathsf T}G^{-1}e\,dz$
proves the first inequality in \eqref{eq:v109-integrable-inverse-gradient}.
The second uses $z_b^2-z_a^2\leq2\Lambda_\circ(b-a)$ and
$(z_b-z_a)^2\leq z_b^2-z_a^2$.
If $z_b=0$, strict increase would contradict $z>0$ in the interior.
\end{proof}

\begin{lemma}[A continuous reduced Hessian at visited grazing]
\label{lem:v109-continuous-schur}
The reduced Hessian has a continuous extension along a fixed
itinerary when any collection of its interior $c_i$ tends to zero,
provided $\tau_i\geq\tau_-$ and \eqref{eq:v109-endpoint-cap} hold.
Its extension still satisfies \eqref{eq:v109-hessian-bounds}.
If at least one interior $c_j$ vanishes, the extended mixed entry
$G_{0m}$ vanishes.
\end{lemma}
\begin{proof}
Let $L=\sum_{i=0}^{m-1}\tau_i^{-1}
(e_i+e_{i+1})(e_i+e_{i+1})^{\mathsf T}$, and set
$C_E=\operatorname{diag}(c_0,c_m)$.
For $m>1$ put
\[
 Z=2I+R C_I^{1/2}L_{II}C_I^{1/2}\geq2I.
\]
Formula \eqref{eq:v109-schur} can be written without inverse incidences:
\begin{equation}\label{eq:v109-continuous-schur}
 \begin{split}
 G={}&RC_E+R^2 C_E L_{EE}C_E\\
 &-R^3 C_E L_{EI}C_I^{1/2}Z^{-1}
                      C_I^{1/2}L_{IE}C_E.
 \end{split}
\end{equation}
Every term is continuous for $c_i\geq0$, since $Z^{-1}$ remains
bounded. The order inequalities pass to this limit. The marked
bound in \eqref{eq:v109-optical-bounds}, together with
$|G_{0m}|=R\omega$, proves the last assertion. There is no interior
extension to make when $m=1$.
\end{proof}

At an endpoint of a fixed-$\theta$ segment the complete finite
trajectory also has a limit. Indeed Lemma~\ref{lem:v109-gauss}
gives the initial angle and $p_0=-z\cos\theta/R$. For each designated
next disk, the entry root of the flight quadratic is continuous up
to a tangency. The flight length cannot collapse to zero, and the
reflection $v\mapsto v-2(v\cdot n)n$ is continuous at tangency.
Induction over the fixed finite itinerary gives the limits of all
visited points, velocities, $\tau_i$ and $c_i$. The third-disk
clearances have limits as well: the finite horizon restricts the
possible nearby disk centres to a finite set, and the minimum
distance to each of them is continuous. This argument takes only
one-sided limits from the actual event; it asserts no continuation
through its boundary.

\begin{theorem}[Trace completion on the original physical sheets]
\label{thm:v109-traces}
On every positive-length original regular pairing segment the
coarea factor and weighted density are
\begin{equation}\label{eq:v109-chi}
 \chi=\det D_q(F,\arg g)=e^{\perp\mathsf T}G e^\perp
                       \geq\lambda_\circ,
 \qquad f^\varepsilon=\frac{\omega W^\varepsilon}{4\pi c_*\chi}.
\end{equation}
They have finite one-sided endpoint limits. The regular-point
variation input \eqref{eq:v109-regular-variation} implies that
$f^\varepsilon$ extends absolutely continuously to the closed
segment and satisfies
\begin{equation}\label{eq:v109-completed-variation}
 |f^\varepsilon(b-)-f^\varepsilon(a+)|
 \leq C(1+\varepsilon^{-1})
       \int_a^b A I_m/z\,du.
\end{equation}
In fact, on each physical sheet,
\begin{equation}\label{eq:v109-holder}
 |f^\varepsilon(v)-f^\varepsilon(u)|
 \leq C(1+\varepsilon^{-1})\varrho^{(m-2)_+}|v-u|^{1/2}
 \quad(a\leq u<v\leq b).
\end{equation}
If a limiting trajectory at either endpoint contains an actually
visited interior grazing collision, its trace at that endpoint is
zero. The original paired minimum on that entire segment is then
zero, for every original selector.
\end{theorem}
\begin{proof}
The determinant identity in \eqref{eq:v109-chi} follows by
expressing $D\arg g$ as $e^{\perp\mathsf T}G/z$.
The preceding endpoint-limit argument and
Lemma~\ref{lem:v109-continuous-schur} give limits for $G$, $\omega$
and $\chi$, the last with a positive lower bound.
Every factor in the finite product defining $W^\varepsilon$ is
continuous in its incidence or clearance argument. Thus
$f^\varepsilon$ has the asserted finite limits, including at a
critical incoming endpoint, where the direction $\theta$ is fixed.

Corollary~\ref{cor:v109-marked} and \eqref{eq:v109-chi} bound
$AI_m$ by $C\varrho^{(m-2)_+}$. The integral in
\eqref{eq:v109-completed-variation} is finite by
Lemma~\ref{lem:v109-gauss}. On compact regular subsegments the flight charts are smooth and
the finite clearance minima are locally Lipschitz, so
$f^\varepsilon$ is locally absolutely continuous. Integrating the
inherited regular-point derivative inequality on these subsegments and then passing to
the proved endpoint limits establishes absolute continuity and
\eqref{eq:v109-completed-variation}. The same argument between $u$
and $v$, using \eqref{eq:v109-integrable-inverse-gradient}, gives
\eqref{eq:v109-holder}.

Finally, if $c_j\to0$ at a visited interior collision, the marked
optical bound gives $\omega\to0$, while $0\leq W^\varepsilon\leq1$
and $\chi\geq\lambda_\circ$. Hence that endpoint trace is zero.
The paired density is its nonnegative minimum with the other trace.
\end{proof}

\begin{remark}[Two different tangency mechanisms]
\label{rem:v109-tangency}
The last assertion concerns a collision in the visited itinerary.
A flight which merely becomes tangent to an \emph{unvisited} third
disk has zero clearance but need not have any $c_j=0$.
If all visited incidences remain positive, then
\eqref{eq:v109-omega} has a positive limit. A zero third-disk
clearance makes $W^\varepsilon=1$ at the limit; consequently the
trace can be positive. Such sheets are not removed. Neither are
nongrazing critical entrances or positive section-boundary traces.
For a nongrazing critical entrance the exact directional limit is
$\omega_*W_*^\varepsilon/(4\pi c_*e^{\perp\mathsf T}G_*e^\perp)$.
No roof-independent value at the critical point is asserted.
\end{remark}
